%% kapa: exam=1 round=35 year=2024 period=1 form=A subject=civil_law subject_ko=민법 seq=34 number=34
%% source: Q-Net 감정평가사 기출문제 (cst00309) artlSeq=5213229 file=raw/1008/5213229/1교시 문제지 원본.pdf
%% conversion: kordoc 4.24.51 -> tools/gen1.py (auto); review=proofread-sonnet; flags=none
\begin{question}{34}
토지전세권에 관한 설명으로 옳은 것을 모두 고른 것은? (다툼이 있으면 판례에 따름)
\begin{framedtext}
ㄱ. 전세권의 존속기간이 만료하면 전세권의 용익물권적 권능은 전세권설정등기의 말소 없이도 당연히 소멸한다. \par
ㄴ. 전세금의 지급은 전세권의 성립요소가 되는 것이므로 기존의 채권으로 전세금지급을 대신할 수 없다. \par
ㄷ. 전세권 존속기간이 시작되기 전에 마친 전세권설정등기도 특별한 사정이 없는 한 유효한 것으로 추정된다. \par
ㄹ. 당사자가 채권담보의 목적으로 전세권을 설정하였으나 설정과 동시에 목적물을 인도하지 않았다면, 장차 전세권자가 목적물을 사용ㆍ수익하기로 하였더라도 그 전세권은 무효이다. \par
\end{framedtext}
\begin{choices}
  \item ㄱ, ㄴ
  \item ㄱ, ㄷ
  \item ㄱ, ㄹ
  \item ㄴ, ㄹ
  \item ㄷ, ㄹ
\end{choices}
\end{question}
\answer{2}
